\section*{Chapter \ref{ch:nw:network-cards-and-kernel-bypass} --- Network Cards and Kernel Bypass}

\begin{solution}{exo:nw:network-cards-and-kernel-bypass:1}
$1.5 \times 10^6 \times 400 \times 10^{-6} = 600$ arrivals for $512 - 40 = 472$ free descriptors: 128 packets dropped.
\end{solution}

\begin{solution}{exo:nw:network-cards-and-kernel-bypass:2}
$1 + \lambda u = 1 + 0.2 \times 20 = 5$ packets an interrupt; interrupts at $200\,000 / 5 = 40\,000$ a second.
\end{solution}

\begin{solution}{exo:nw:network-cards-and-kernel-bypass:3}
$15.4 / 16.2$: 95\%. Almost all of it is the wake-up and scheduling of the sleeping reader, not the network stack.
\end{solution}

\begin{solution}{exo:nw:network-cards-and-kernel-bypass:4}
RSS hashes flows over rings to balance cores; a feed handler wants all of a line's groups on the one ring its pinned core polls, and no
other traffic there. \omterm{def:nw:network-cards-and-kernel-bypass:rss}{Flow steering} pins the venue's flows to the chosen ring and keeps the rest of the host's traffic on others.
\end{solution}

\begin{solution}{exo:nw:network-cards-and-kernel-bypass:5}
$1/(2 + 0.3) = 0.435$ packets a microsecond: 435\,000 a second. With 16 packets an interrupt, $1/(2/16 + 0.3) = 2.35$ million a second.
\end{solution}

\begin{solution}{exo:nw:network-cards-and-kernel-bypass:6}
The poller never sleeps, so the wake-up (most of the blocking median) disappears. But on shared cores the scheduler sometimes runs
something else on the poller's core for milliseconds, and every datagram arriving meanwhile waits. On a production host the polling
thread gets an isolated core with its interrupts and timer tick moved away (One Quant Book~13, chapter~13), on the card's socket.
\end{solution}

\begin{solution}{exo:nw:network-cards-and-kernel-bypass:7}
128 descriptors: with batches of 32, a ring of 64 still drops 36 packets of the fixture, a ring of 128 drops none.
\end{solution}

\begin{solution}{exo:nw:network-cards-and-kernel-bypass:8}
Loss is rare, not impossible (a flapping optic, a switch buffer during a burst, a failover): without retransmission one lost segment
leaves the session stuck or silently missing an order, and the venue's session layer will disconnect or cancel. The two microseconds
came from removing the kernel, not from removing reliability; use a user-space or hardware TCP that implements it.
\end{solution}

\begin{solution}{pb:nw:network-cards-and-kernel-bypass:1}
\begin{enumerate}
\item $0.2 \times (2 + 0.3) = 0.46$: 46\% of a core.
\item $3 + 2 + 0.3 = \qty{5.3}{\micro\second}$.
\item Queueing: an interrupt that arrives while the previous one is still being handled waits for it.
\item $1/2.3$ per microsecond: 435\,000 packets a second.
\item $1 + 0.2 \times 10 = 3$.
\item $200\,000 / 3 \approx 66\,667$ a second.
\item $66\,667 \times 2 \times 10^{-6} + 0.2 \times 0.3 = 0.133 + 0.06$: 19.3\%.
\item From the timer: the first packet of a batch waits \qty{10}{\micro\second} for it, later ones less, before the \qty{5.3}{\micro\second}
  of handling.
\item Half a loop, \qty{0.1}{\micro\second}.
\item $0.1 + 0.3 = \qty{0.4}{\micro\second}$; the model's 0.41 adds the rare packet that arrives while the previous one is processed.
\item The whole core.
\item Isolated, never descheduled, on the card's socket, with nothing else to do.
\item \textbf{Named result.} At 200\,000 packets a second: a \qty{10}{\micro\second} timer gives \qty{12.4}{\micro\second} of mean latency
  for 19.3\% of a core; polling gives \qty{0.41}{\micro\second} for all of it; interrupts alone would take a whole core at 435\,000 packets
  a second.
\item About \qty{12}{\micro\second} (12.4 against 0.41).
\item Where $\lambda(2/(1 + 10\lambda) + 0.3) = 0.25$: about 324\,000 packets a second.
\item The hot path is worth a core; the rest of the host's traffic is not, and coalescing keeps it from interrupting the machine all day.
\item Interrupt cost and wake-up with a profiler and hardware timestamps, the card's coalescing behaviour under the real feed, and the
  distribution of pauses of the polling core.
\item $1\,024 / 200\,000$ seconds: \qty{5.1}{\milli\second} of pause before the ring overflows.
\item Microseconds of wake-up per batch, and a tail as long as the scheduler's longest decision.
\item Raw ring access for the multicast feed, polled on an isolated core: the feed needs no TCP, and raw access removes the most.
\end{enumerate}
\end{solution}

\begin{solution}{iq:nw:network-cards-and-kernel-bypass:1}
The card writes the frame by DMA into a buffer named by the next receive descriptor and marks it done; it raises an interrupt (perhaps
after coalescing); the kernel schedules NAPI processing in a software interrupt, which reads the ring, builds the socket buffer, passes
it up IP and UDP, queues it on the socket and wakes the reader; \texttt{recv} copies it into the program's buffer.
\iqlookfor{DMA, the ring, the interrupt and softirq, the socket queue, the copy and the wake-up.}
\end{solution}

\begin{solution}{iq:nw:network-cards-and-kernel-bypass:2}
The card delays the receive interrupt by a time or a packet count so that one interrupt serves several packets: less CPU, more latency.
On a market-data host the hot rings are polled, so their coalescing does not matter; for the others choose moderate values and turn
adaptive coalescing off where predictability matters.
\iqlookfor{the latency-CPU trade-off and that the hot path avoids it by polling.}
\end{solution}

\begin{solution}{iq:nw:network-cards-and-kernel-bypass:3}
Running the data path in user space: raw ring access (vendor interfaces, DPDK), user-space stacks behind the sockets API, or AF\_XDP.
It removes interrupts, system calls, copies and wake-ups. You give up the kernel's tools, isolation and TCP (unless the stack brings
its own), and you spend a core per polling thread.
\iqlookfor{the families and the costs, not only the speed.}
\end{solution}

\begin{solution}{iq:nw:network-cards-and-kernel-bypass:4}
Whether the polling core is really isolated (other tasks, interrupts, timer ticks, kernel threads on it), frequency and idle states,
page faults and allocation on the path, the ring size against the longest pause, and whether the tail coincides with bursts. Measure
with hardware timestamps against the software ones.
\iqlookfor{scheduling interference first, then the ring.}
\end{solution}

\begin{solution}{iq:nw:network-cards-and-kernel-bypass:5}
Average load says nothing about pauses: while the host is not processing (a burst, a descheduling, a batched refill), the card consumes
free descriptors at the arrival rate, and a pause longer than free descriptors divided by that rate drops packets. Size the ring for the
burst rate times the longest pause plus the refill batch.
\iqlookfor{pause times rate, and the descriptors held by batching.}
\end{solution}

\begin{solution}{iq:nw:network-cards-and-kernel-bypass:6}
The kernel's timestamp is taken after the card, its ring, the interrupt and the driver; the card's is taken at the port. Their
difference is the receive path the firm is trying to tune, and only the card's timestamp compares with the timestamps of switches and
taps (chapter~5).
\iqlookfor{where on the path each timestamp is taken.}
\end{solution}
